% Kreis – bank/kreis/ (zusammenbau v0.4, Heft)
\einheitenkopf[t6]{Kreis}
\setcounter{aufgabe}{43}

% Anlauf (ohne Prüfkennung)
\begin{aufgabe}{Fläche – Anlauf}
\begin{teile}
\teil Wie viel Stoff braucht man für eine runde Tischdecke? Welche Formel brauchst du? Kreuze an.\\ \kreuz{$u = \pi \cdot d$}\\ \kreuz{$A = \pi \cdot r^2$}
\teil Kreis mit $r = 13$ cm – wie groß ist die Fläche? Rechne mit der $\pi$-Taste und runde auf eine Stelle. A ≈ \leerfeld[cm²]
\teil Kreis mit $d = 16$ cm – wie groß ist die Fläche? Rechne mit der $\pi$-Taste und runde auf eine Stelle. A ≈ \leerfeld[cm²]
\end{teile}
\end{aufgabe}

\begin{aufgabe}{Fläche – Prüfungsaufgaben}
\begin{teile}
\teil Ein runder Sandkasten hat den Durchmesser $2{,}36$ m. Berechne seinen Flächeninhalt. Runde auf zwei Stellen \hfill (P10 2016 OS). A ≈ \leerfeld
\teil Ein rundes Trampolin hat den Durchmesser $3{,}66$ m. Berechne den Flächeninhalt der Sprungfläche. Runde auf zwei Stellen \hfill (P10 2016 OS). A ≈ \leerfeld
\teil Ein Sandhaufen hat die Form eines Kegels mit dem Radius $45$ cm und der Höhe $60$ cm. Berechne den Flächeninhalt der Grundfläche. Runde auf ganze cm² \hfill (P10 2024 OS). A ≈ \leerfeld
\teil Ein Zelt hat die Form eines Kegels mit dem Radius $1{,}65$ m und der Höhe $2{,}4$ m. Berechne den Flächeninhalt des Zeltbodens. Runde auf zwei Stellen \hfill (P10 2024 OS). A ≈ \leerfeld
\end{teile}
\end{aufgabe}

% Anlauf (ohne Prüfkennung)
\begin{aufgabe}{Kreisteile – Anlauf}
\begin{teile}
\teil Welcher Teil des Kreises ist grau? Kreuze an.\\ \kreuz{$\frac{1}{2}$}\\ \kreuz{$\frac{1}{3}$}\\ \kreuz{$\frac{1}{4}$}

\kreissektor{180}{}
\teil Ein Kreisausschnitt hat den Mittelpunktswinkel $\alpha = 90^\circ$. Welcher Teil des Kreises ist das? Gib den Anteil als Bruch an. \leerfeld
\teil Ein Kreisausschnitt hat den Mittelpunktswinkel $\alpha = 100^\circ$. Wie viel Prozent des Kreises sind das? Runde auf eine Stelle. ≈ \leerfeld[\%]
\end{teile}
\end{aufgabe}

\begin{aufgabe}{Kreisteile – Prüfungsaufgaben}
\begin{teile}
\teil Im Bild ist ein Kreisausschnitt grau gefärbt. Sein Mittelpunktswinkel beträgt $130^\circ$, der Winkel des weißen Teils heißt $\alpha$. Wie viel Prozent der Kreisfläche sind grau? Runde auf eine Stelle \hfill (P10 2025 OS).

\kreissektor{130}{$130^\circ$}
≈ \leerfeld[\%]
\teil Ein rundes Beet ist geteilt: Auf dem grauen Ausschnitt im Bild mit dem Mittelpunktswinkel $155^\circ$ wachsen Rosen, auf dem weißen Teil mit dem Winkel $\alpha$ Gras. Wie viel Prozent des Beets tragen Rosen? Runde auf eine Stelle \hfill (P10 2025 OS).

\kreissektor{155}{$155^\circ$}
≈ \leerfeld[\%]
\end{teile}
\end{aufgabe}
